\section{Задание 3. Код Хэмминга (7, 4)}

\subsection{Постановка задачи}
В данном задании требуется освоить корректирующий код Хэмминга с параметрами $(7, 4)$, который позволяет обнаруживать и исправлять одиночные ошибки в 7-битных кодовых словах при передаче 4 информационных бит. Необходимо составить матрицы генератора $G$ и проверочную $H$, закодировать интересное слово из 4 букв русского алфавита, сымитировать вредоносное вмешательство (замену 1, 2, 3 и 4 бит) и восстановить исходное сообщение.

\subsection{Алфавит и двоичное представление}
Для работы выбран русский алфавит из 32 букв (без буквы \texttt{Ё}). Каждой букве сопоставлен пятибитовый двоичный номер от \texttt{00000} до \texttt{11111}. Таблица соответствия приведена в таблице \ref{tab:alphabet}.

\begin{center}
\small
\begin{tabular}{|c|c|c|c|c|c|c|c|}
\hline
\textbf{А} & \textbf{Б} & \textbf{В} & \textbf{Г} & \textbf{Д} & \textbf{Е} & \textbf{Ж} & \textbf{З} \\
\hline
00000 & 00001 & 00010 & 00011 & 00100 & 00101 & 00110 & 00111 \\
\hline
\textbf{И} & \textbf{Й} & \textbf{К} & \textbf{Л} & \textbf{М} & \textbf{Н} & \textbf{О} & \textbf{П} \\
\hline
01000 & 01001 & 01010 & 01011 & 01100 & 01101 & 01110 & 01111 \\
\hline
\textbf{Р} & \textbf{С} & \textbf{Т} & \textbf{У} & \textbf{Ф} & \textbf{Х} & \textbf{Ц} & \textbf{Ч} \\
\hline
10000 & 10001 & 10010 & 10011 & 10100 & 10101 & 10110 & 10111 \\
\hline
\textbf{Ш} & \textbf{Щ} & \textbf{Ъ} & \textbf{Ы} & \textbf{Ь} & \textbf{Э} & \textbf{Ю} & \textbf{Я} \\
\hline
11000 & 11001 & 11010 & 11011 & 11100 & 11101 & 11110 & 11111 \\
\hline
\end{tabular}
\end{center}
\label{tab:alphabet}

Исходное слово из 4 букв: \textbf{ТЕСТ}. Каждая буква преобразуется в 5-битный код:

\begin{align}
    \text{Т} \rightarrow 10010, \\
    \text{Е} \rightarrow 00101, \\
    \text{С} \rightarrow 10001, \\
    \text{Т} \rightarrow 10010.
\end{align}

Полный двоичный код слова (20 символов): \texttt{10010001011000110010}.

\subsection{Матрицы генератора $G$ и проверочная $H$}
Для кода Хэмминга $(7, 4)$ выбраны стандартные матрицы. Матрица генератора $G$ размером $4 \times 7$ составляется как $G = [I_4 \;|\; P]$, где $P$ — матрица проверки чётности $4 \times 3$:

\[
G = \begin{pmatrix}
1 & 0 & 0 & 0 & 1 & 1 & 0 \\
0 & 1 & 0 & 0 & 1 & 0 & 1 \\
0 & 0 & 1 & 0 & 0 & 1 & 1 \\
0 & 0 & 0 & 1 & 1 & 1 & 1
\end{pmatrix}
\]

Проверочная матрица $H$ размером $3 \times 7$ составляется как $H = [P^T \;|\; I_3]$:

\[
H = \begin{pmatrix}
1 & 1 & 0 & 1 & 1 & 0 & 0 \\
1 & 0 & 1 & 1 & 0 & 1 & 0 \\
0 & 1 & 1 & 1 & 0 & 0 & 1
\end{pmatrix}
\]

Проверка корректности: $G \cdot H^T \equiv 0 \pmod{2}$, что подтверждает линейную зависимость строк и правильное построение кода.

\begin{center}
\fbox{\parbox{0.9\textwidth}{
\centering
\textbf{Матрица генератора $G$ (4$\times$7)} \\
$G = \begin{pmatrix}
1 & 0 & 0 & 0 & 1 & 1 & 0 \\
0 & 1 & 0 & 0 & 1 & 0 & 1 \\
0 & 0 & 1 & 0 & 0 & 1 & 1 \\
0 & 0 & 0 & 1 & 1 & 1 & 1
\end{pmatrix}$ \\
\vspace{4pt}
\textbf{Проверочная матрица $H$ (3$\times$7)} \\
$H = \begin{pmatrix}
1 & 1 & 0 & 1 & 1 & 0 & 0 \\
1 & 0 & 1 & 1 & 0 & 1 & 0 \\
0 & 1 & 1 & 1 & 0 & 0 & 1
\end{pmatrix}$
}}
\end{center}

\subsection{Процесс кодирования}
Полный 20-битный код разбивается на 5 блоков по 4 информационных бита. Каждый блок умножается на $G$ в поле $\mathbb{F}_2$ (по модулю 2), что даёт 7-битное кодовое слово. Результаты кодирования приведены в таблице \ref{tab:encode}.

\begin{table}[H]
    \centering
    \small
    \caption{Кодирование блоков с помощью матрицы $G$}
    \label{tab:encode}
    \begin{tabular}{|c|c|c|}
        \toprule
        Блок (4 бита) & Код (7 бит) & Десятичное значение \\
        \midrule
        1001 & 1001001 & 73 \\
        0001 & 0001111 & 15 \\
        0110 & 0110110 & 54 \\
        0011 & 0011100 & 28 \\
        0010 & 0010011 & 19 \\
        \bottomrule
    \end{tabular}
\end{table}

Полный закодированный код (35 бит) составляет:
\texttt{1001001 0001111 0110110 0011100 0010011}.

Процесс кодирования визуализирован на рисунке \ref{fig:flow}.

\begin{center}
\fbox{\parbox{0.95\textwidth}{
\small
\textbf{Схема кодирования} \quad ТЕСТ $\rightarrow$ 10010 00101 10001 10010 $\rightarrow$ 5 блоков по 4 бита \quad $\rightarrow$ \quad 1001001 \; 0001111 \; 0110110 \; 0011100 \; 0010011 \quad $\rightarrow$ \quad 35 бит
}}
\end{center}
\label{fig:flow}

\subsection{Симуляция вредоносного вмешательства}
Для демонстрации возможностей кода последовательно испорчены 1, 2, 3 и 4 бита в различных позициях 35-битного кода. Искажённые сообщения с выделенными изменёнными битами приведены в таблице \ref{tab:corrupt}.

\begin{table}[H]
    \centering
    \small
    \caption{Искажённые сообщения с выделенными ошибочными битами}
    \label{tab:corrupt}
    \begin{tabular}{|l|l|}
        \toprule
        Вмешательство & Искажённый код (выделены ошибки) \\
        \midrule
        1 бит (позиция 5) & \texttt{10010[1]1 0001111 0110110 0011100 0010011} \\
        2 бита (3, 15) & \texttt{100[0]001 0001111 011[0]110 0011100 00[1]0011} \\
        3 бита (8, 20, 30) & \texttt{1001001 0[1]01111 011011[1] 0011100 00[0]0011} \\
        4 бита (1, 10, 25, 32) & \texttt{1[1]01001 000[0]111 0110110 0011[0]00 0010[1]11} \\
        \bottomrule
    \end{tabular}
\end{table}

\subsection{Декодирование и исправление ошибок}
Декодирование выполняется с помощью матрицы $H$. Для каждого 7-битного блока вычисляется синдром $s = H \cdot c^T \pmod{2}$. Синдром однозначно указывает позицию ошибочного бита (если он один в блоке). Таблица соответствия синдромов приведена на рисунке \ref{fig:syndrome}.

\begin{center}
\fbox{\parbox{0.92\textwidth}{
\small
\textbf{Процесс исправления ошибок с матрицей $H$} \\
1. Вычисляем синдром: $s = H \cdot c^T \pmod{2}$ \\
2. Синдром указывает позицию ошибочного бита \\
3. Инвертируем найденный бит \\
4. Извлекаем информационные биты (первые 4) \\
\vspace{4pt}
\textbf{Таблица синдромов} \\
\begin{tabular}{|c|c|}
\hline
\textbf{Синдром} & \textbf{Позиция ошибки} \\
\hline
000 & нет ошибки \\
110 & бит 1 \\
101 & бит 2 \\
011 & бит 3 \\
111 & бит 4 \\
100 & бит 5 (p1) \\
010 & бит 6 (p2) \\
001 & бит 7 (p3) \\
\hline
\end{tabular}
}}
\end{center}
\label{fig:syndrome}

Для каждого случая вмешательства выполнена декодировка. Результаты приведены в таблице \ref{tab:decode}. При 1--3 ошибках в разных блоках код успешно исправляет все ошибки и восстанавливает исходное слово \textbf{ТЕСТ}. При 4 ошибках, распределённых по разным блокам, также достигается полное восстановление, что демонстрирует высокую устойчивость кода при правильном распределении ошибок.

\begin{table}[H]
    \centering
    \caption{Результаты декодирования после исправления ошибок}
    \label{tab:decode}
    \begin{tabular}{|l|l|l|}
        \toprule
        Вмешательство & Декодированное слово & Верно? \\
        \midrule
        1 бит & ТЕСТ & ДА \\
        2 бита & ТЕСТ & ДА \\
        3 бита & ТЕСТ & ДА \\
        4 бита & ТЕСТ & ДА \\
        \bottomrule
    \end{tabular}
\end{table}

\subsection{Расчёты исправления ошибок}
Приведём пример расчёта для блока с одной ошибкой. Пусть кодовый блок $c = [1, 0, 1, 1, 0, 0, 1]$ (ошибка в бите 2). Вычисляем синдром:

\[
s = H \cdot c^T = \begin{pmatrix}
1 & 1 & 0 & 1 & 1 & 0 & 0 \\
1 & 0 & 1 & 1 & 0 & 1 & 0 \\
0 & 1 & 1 & 1 & 0 & 0 & 1
\end{pmatrix}
\begin{pmatrix}
1 \\ 0 \\ 1 \\ 1 \\ 0 \\ 0 \\ 1
\end{pmatrix}
= \begin{pmatrix}
1 \oplus 0 \oplus 1 \oplus 0 \oplus 0 \\ 1 \oplus 1 \oplus 1 \oplus 0 \oplus 0 \\ 0 \oplus 0 \oplus 1 \oplus 0 \oplus 1
\end{pmatrix}
= \begin{pmatrix}
1 \\ 0 \\ 1
\end{pmatrix} = (101)_2.
\]

Синдром $(101)_2$ соответствует ошибке в бите 2. После инвертирования второго бита и извлечения первых 4 информационных бит получаем исходные данные блока.

\subsection{Выводы по заданию 3}
В ходе выполнения задания были составлены матрицы $G$ и $H$ для кода Хэмминга $(7, 4)$, выполнено кодирование и декодирование 20-битного сообщения, сымитировано вредоносное вмешательство с заменой 1--4 бит. Показано, что код позволяет обнаруживать и исправлять одиночные ошибки в каждом 7-битном блоке. При распределении ошибок по разным блокам код успешно восстанавливает исходное слово даже при 4 ошибках. Это демонстрирует эффективность корректирующих кодов для защиты информации при передаче.
